% GENERATED FILE. Edit canonical lessons, then run npm run generate:books.
% canonical-source-hash: fnv1a64:5e147e981d1a9398
% canonical-lessons: JA-C294-mikoshi, JA-C294-katsugu, JA-C294-gyouretsu, JA-C294-moriagaru, JA-C294-bon-odori

\chapter{Carrying the Shrine}
\label{ch:ja-carrying-the-shrine}
\hlchaptermodality{294}

\begin{chapteropening}
\textbf{By the end of this chapter:} \emph{I can say what the festival crowd carried and danced.}

Complete the last lesson of chapter 294: I can say what the festival crowd carried and danced.
\end{chapteropening}

\section[\texorpdfstring{\ja{みこし} (mikoshi)}{みこし (mikoshi)}]{\ja{みこし} (mikoshi) — a portable shrine}

\label{lesson:JA-C294-mikoshi}



\begin{quote}
Before the new one: say the Japanese for the night sky, then the Japanese for to bustle.
\end{quote}

\subsection*{You'll want to know: \ja{みこし}}
\textbf{\ja{みこし}} — \emph{mikoshi} — \textquotedblleft{}a portable shrine\textquotedblright{}.

\begin{quote}
\textbf{\ja{まつりで} \ja{みこしが} \ja{まちを} \ja{とおります。}} — \emph{matsuri de mikoshi ga machi o tōrimasu} — \textquotedblleft{}At the festival a portable shrine goes through the town.\textquotedblright{}
\end{quote}

\begin{cousinweb}
\textbf{\ja{みこし}} (\emph{mikoshi}) is a native Japanese word (wago), inherited rather than borrowed from Chinese; its kanji was attached later as a kun'yomi (JMdict).
\end{cousinweb}

\begin{grammarlens}[title={What you've built}]
One more word for telling what happened.
\end{grammarlens}

\subsection*{Your turn}
\begin{itemize}
\raggedright
  \item \emph{Say it:} \emph{mikoshi}
  \item \emph{Say it:} \emph{mikoshi}, once more
  \item \emph{Say it:} \emph{matsuri de mikoshi ga machi o tōrimasu} — At the festival a portable shrine goes through the town.
  \item \emph{Recall:} say the Japanese for the night sky, then the Japanese for to bustle, then say \emph{mikoshi} again
\end{itemize}

\begin{tcolorbox}[breakable,colback=teal!4,colframe=teal!35!black,title={Before you move on}]
What does \ja{みこし} mean? (\textquotedblleft{}A portable shrine\textquotedblright{}.) Say it once more.
\end{tcolorbox}
\section[\texorpdfstring{\ja{かつぐ} (katsugu)}{かつぐ (katsugu)}]{\ja{かつぐ} (katsugu) — to carry on one's shoulders}

\label{lesson:JA-C294-katsugu}



\begin{quote}
Before the new one: say the Japanese for to bustle, then the Japanese for a portable shrine.
\end{quote}

\subsection*{You'll want to know: \ja{かつぐ}}
\textbf{\ja{かつぐ}} — \emph{katsugu} — \textquotedblleft{}to carry on one's shoulders\textquotedblright{}.

\begin{quote}
\textbf{\ja{わかい} \ja{おとこが} \ja{みこしを} \ja{かついでいます。}} — \emph{wakai otoko ga mikoshi o katsuide imasu} — \textquotedblleft{}Young men are carrying the shrine on their shoulders.\textquotedblright{}
\end{quote}

Polite \textbf{\ja{かつぎます}}, \emph{katsugimasu}; a godan verb. Past (polite): \textbf{\ja{あにが} \ja{みこしを} \ja{かつぎました。}}, \emph{ani ga mikoshi o katsugimashita}.

\begin{cousinweb}
\textbf{\ja{かつぐ}} (\emph{katsugu}) is a native Japanese word (wago), inherited rather than borrowed from Chinese; its kanji was attached later as a kun'yomi (JMdict).
\end{cousinweb}

\begin{grammarlens}[title={What you've built}]
One more word for telling what happened.
\end{grammarlens}

\subsection*{Your turn}
\begin{itemize}
\raggedright
  \item \emph{Say it:} \emph{katsugu}
  \item \emph{Say it:} \emph{katsugu}, once more
  \item \emph{Say it:} \emph{wakai otoko ga mikoshi o katsuide imasu} — Young men are carrying the shrine on their shoulders.
  \item \emph{Recall:} say the Japanese for to bustle, then the Japanese for a portable shrine, then say \emph{katsugu} again
\end{itemize}

\begin{tcolorbox}[breakable,colback=teal!4,colframe=teal!35!black,title={Before you move on}]
What does \ja{かつぐ} mean? (\textquotedblleft{}To carry on one's shoulders\textquotedblright{}.) Say it once more.
\end{tcolorbox}
\section[\texorpdfstring{\ja{ぎょうれつ} (gyōretsu)}{ぎょうれつ (gyōretsu)}]{\ja{ぎょうれつ} (gyōretsu) — a procession; a queue}

\label{lesson:JA-C294-gyouretsu}



\begin{quote}
Before the new one: say the Japanese for a portable shrine, then the Japanese for to carry on one's shoulders.
\end{quote}

\subsection*{You'll want to know: \ja{ぎょうれつ}}
\textbf{\ja{ぎょうれつ}} — \emph{gyōretsu} — \textquotedblleft{}a procession; a queue\textquotedblright{}.

\begin{quote}
\textbf{\ja{ぎょうれつは} \ja{ゆっくり} \ja{すすみました。}} — \emph{gyōretsu wa yukkuri susumimashita} — \textquotedblleft{}The procession moved slowly along.\textquotedblright{}
\end{quote}

\begin{cousinweb}
\textbf{\ja{ぎょうれつ}} (\emph{gyōretsu}) is Sino-Japanese (kango): two Chinese characters read with their on'yomi, readings borrowed from Chinese. They are \emph{gyō}, \textquotedblleft{}line\textquotedblright{} (Middle Chinese \emph{hang}) and \emph{retsu}, \textquotedblleft{}row\textquotedblright{} (\emph{ljet}) (KANJIDIC2; Middle Chinese in Baxter's transcription).
\end{cousinweb}

\begin{grammarlens}[title={What you've built}]
One more word for telling what happened.
\end{grammarlens}

\subsection*{Your turn}
\begin{itemize}
\raggedright
  \item \emph{Say it:} \emph{gyōretsu}
  \item \emph{Say it:} \emph{gyōretsu}, once more
  \item \emph{Say it:} \emph{gyōretsu wa yukkuri susumimashita} — The procession moved slowly along.
  \item \emph{Recall:} say the Japanese for a portable shrine, then the Japanese for to carry on one's shoulders, then say \emph{gyōretsu} again
\end{itemize}

\begin{tcolorbox}[breakable,colback=teal!4,colframe=teal!35!black,title={Before you move on}]
What does \ja{ぎょうれつ} mean? (\textquotedblleft{}A procession; a queue\textquotedblright{}.) Say it once more.
\end{tcolorbox}
\section[\texorpdfstring{\ja{もりあがる} (moriagaru)}{もりあがる (moriagaru)}]{\ja{もりあがる} (moriagaru) — to liven up; to get going}

\label{lesson:JA-C294-moriagaru}



\begin{quote}
Before the new one: say the Japanese for to carry on one's shoulders, then the Japanese for a procession.
\end{quote}

\subsection*{You'll want to know: \ja{もりあがる}}
\textbf{\ja{もりあがる}} — \emph{moriagaru} — \textquotedblleft{}to liven up; to get going\textquotedblright{}.

\begin{quote}
\textbf{\ja{ゆうがたに} \ja{なって、まつりは} \ja{もりあがりました。}} — \emph{yūgata ni natte, matsuri wa moriagarimashita} — \textquotedblleft{}As evening came on, the festival really got going.\textquotedblright{}
\end{quote}

Polite \textbf{\ja{もりあがります}}, \emph{moriagarimasu}; a godan verb. Past (polite): \textbf{\ja{まつりが} \ja{もりあがりました。}}, \emph{matsuri ga moriagarimashita}.

\begin{cousinweb}
\textbf{\ja{もりあがる}} (\emph{moriagaru}) is a native Japanese word (wago), inherited rather than borrowed from Chinese, built from \emph{mori}, \textquotedblleft{}heaping\textquotedblright{} (from \emph{moru}), and \emph{agaru}, \textquotedblleft{}to rise\textquotedblright{} (JMdict, KANJIDIC2).
\end{cousinweb}

\begin{grammarlens}[title={What you've built}]
One more word for telling what happened.
\end{grammarlens}

\subsection*{Your turn}
\begin{itemize}
\raggedright
  \item \emph{Say it:} \emph{moriagaru}
  \item \emph{Say it:} \emph{moriagaru}, once more
  \item \emph{Say it:} \emph{yūgata ni natte, matsuri wa moriagarimashita} — As evening came on, the festival really got going.
  \item \emph{Recall:} say the Japanese for to carry on one's shoulders, then the Japanese for a procession, then say \emph{moriagaru} again
\end{itemize}

\begin{tcolorbox}[breakable,colback=teal!4,colframe=teal!35!black,title={Before you move on}]
What does \ja{もりあがる} mean? (\textquotedblleft{}To liven up; to get going\textquotedblright{}.) Say it once more.
\end{tcolorbox}
\section[\texorpdfstring{\ja{ぼんおどり} (bon'odori)}{ぼんおどり (bon'odori)}]{\ja{ぼんおどり} (bon'odori) — a Bon festival dance}

\label{lesson:JA-C294-bon-odori}



\begin{quote}
Before the new one: say the Japanese for a procession, then the Japanese for to liven up.
\end{quote}

\subsection*{You'll want to know: \ja{ぼんおどり}}
\textbf{\ja{ぼんおどり}} — \emph{bon'odori} — \textquotedblleft{}a Bon festival dance\textquotedblright{}.

\begin{quote}
\textbf{\ja{なつには} \ja{日本中}\ja{で} \ja{ぼんおどりが} \ja{あります。}} — \emph{natsu ni wa nihonjū de bon'odori ga arimasu} — \textquotedblleft{}In summer there is festival dancing all over Japan.\textquotedblright{}
\end{quote}

\begin{cousinweb}
\textbf{\ja{ぼんおどり}} (\emph{bon'odori}) is a mixed word, native and Sino-Japanese: \emph{bon}, \textquotedblleft{}lantern festival\textquotedblright{}, is a Chinese character read with its on'yomi, a reading borrowed from Chinese (Middle Chinese \emph{bwon}), and \emph{odori}, \textquotedblleft{}dance\textquotedblright{}, is native Japanese (JMdict, KANJIDIC2; Middle Chinese in Baxter's transcription).
\end{cousinweb}

\begin{grammarlens}[title={What you've built}]
One more word for telling what happened.
\end{grammarlens}

\subsection*{Your turn}
\begin{itemize}
\raggedright
  \item \emph{Say it:} \emph{bon'odori}
  \item \emph{Say it:} \emph{bon'odori}, once more
  \item \emph{Say it:} \emph{natsu ni wa nihonjū de bon'odori ga arimasu} — In summer there is festival dancing all over Japan.
  \item \emph{Recall:} say the Japanese for a procession, then the Japanese for to liven up, then say \emph{bon'odori} again
\end{itemize}

\begin{tcolorbox}[breakable,colback=teal!4,colframe=teal!35!black,title={Before you move on}]
What does \ja{ぼんおどり} mean? (\textquotedblleft{}A Bon festival dance\textquotedblright{}.) Say it once more.
\end{tcolorbox}
